\documentclass[10pt,a4paper]{article}
\usepackage[utf8]{inputenc}
\usepackage[a4paper,left=3cm,right=2cm]{geometry}
\usepackage{amsmath}
\usepackage[thmmarks, amsmath, thref]{ntheorem}
\usepackage{amsfonts}
\usepackage{amssymb}
\usepackage{fancyhdr}
\pagestyle{fancy}
\usepackage{anysize}
\usepackage{graphicx}
\usepackage{float}
\usepackage{hyperref}
\usepackage{enumitem}
\newtheorem{theorem}{Theorem}
\theoremsymbol{\ensuremath{\blacksquare}}
\newtheorem*{solution}{Solución:}

\usepackage{xcolor}

% Margenes y formalidades

\textwidth 6.25in % Margen Derecho y ancho del texto
\textheight 8.5in
\oddsidemargin 0in % Margen Izquierdo
\headheight 0in % Largo del texto
\leftmargin 1in
\parindent 0em % Salto entre la indentación (Sangría)https://www.overleaf.com/4519127669sdpptkhvsbfn
\topmargin -0.5in % Margen Superior
\parskip 0.5ex % Salto entre parrafos
\baselineskip 1pt

%\lhead{\includegraphics[scale=0.2]{logo.PNG}} % texto izquierda de la cabecera
%\chead{}                                      % texto centro de la cabecera
\rhead{Universidad Técnica Federico Santa María\\Ingeniería Comercial} % número de página a la derecha
\renewcommand{\headrulewidth}{0.4pt} % grosor de la línea de la cabecera
\renewcommand{\footrulewidth}{0.4pt} % grosor de la línea del pie

\begin{document}
\vspace{5cm}


\begin{center}
\begin{large}
\textbf{Ejercicio heterocedasticidad\\}
\end{large}\\
\end{center}

\begin{center}
\begin{large}
\textbf{Econometría ICS-294}
\end{large}\\

\end{center}


Usando los datos del archivo GPA3 se estimó la siguiente ecuación para los estudiantes
que en otoño están en el segundo semestre:



{\scriptsize
\begin{equation*}
\begin{array}{ll}
\widehat{trmgpa} = & -2,12  + 0,193 cumgpa + 0,0014 tothrs + 0,0018 sat - 0,0039 hsperrc + 0,351 female - 0,157 season
 \\
& \quad (0,55)  \quad \quad (0,064)  \quad \quad \quad (0,0012) \quad \quad \quad (0,0002) \quad \quad \quad (0,0018) \quad \quad \quad \quad (0,085) \quad \quad \quad (0,098) \\ 
& \quad [0,55]  \quad \quad [0,074] \quad \quad \quad \quad [0,0012] \quad \quad \quad [0,0002] \quad \quad \quad [0,0019] \quad \quad \quad \quad \quad[0,079] \quad \quad \quad [0,07] \\
& n = \text{269} \quad R^2 = \text{0,465}
\end{array}
\end{equation*}}



Aquí $trmgpa$ es el promedio general de calificaciones (GPA) del semestre actual, $crsgpa$ es un
promedio ponderado de calificaciones de los cursos que se están tomando, $cumgpa$ es el promedio
general de calificaciones antes del semestre presente, $tothrs$ es el total de horas crédito antes
de este semestre, $sat$ es la puntuación en la prueba SAT de admisión a la universidad, $hsperc$
es el percentil que ocupó el alumno entre los graduados del bachillerato, $female$ es una variable
binaria para el género femenino, y $season$ es una variable binaria que es igual a uno si el estudiante practica deporte en otoño. Entre paréntesis y entre corchetes se dan, respectivamente,
los errores estándar usuales y los errores estándar robustos a la heterocedasticidad.


\begin{enumerate}
    \item ¿Tienen las variables  $cumgpa$ y $tothrs$ los efectos estimados esperados? ¿Cuáles
de estas variables son estadísticamente significativas al nivel de 5\%? ¿Importa qué error
estándar se use?

{\color{blue}



% 1. **\(crsgpa\):** El coeficiente es \(0.9\), positivo, lo que indica que un mayor promedio ponderado de calificaciones en los cursos actuales está asociado con un aumento en el \(trmgpa\). Este resultado es consistente con lo esperado, ya que los estudiantes con mejor desempeño en sus cursos actuales suelen tener un mejor GPA del semestre.
\begin{itemize}
    \item \(cumgpa\): El coeficiente es \(0.193\), positivo, lo cual sugiere que un mejor promedio acumulado previo está asociado con un mayor \(trmgpa\). Esto es coherente con las expectativas, ya que los estudiantes con buenos historiales académicos suelen mantener un buen desempeño.
    \item \(tothrs\): El coeficiente es \(0.0014\), positivo pero de magnitud pequeña. Esto indica que más horas crédito acumuladas están levemente asociadas con un mayor \(trmgpa\), lo cual es razonable, ya que la experiencia acumulada puede influir positivamente en el desempeño académico.
% \item \(sat\): Efecto positivo sobre el promedio general, es el efecto esperado ya que mejor puntaje de la prueba de admisión puede influir positivamente en los resultados actuales. 
% \item \(hsperc \): Efecto negativo sobre el resultado, que es lo esperado dado a que mejor ranking implica un valor numerico más bajo  en hsperc. Mientras mejor el ranking de graduación en el bachillerato se esperan  mejores resultados actuales.
\end{itemize}



Para evaluar la significancia estadística de estas variables, calculamos los valores \(t\) como:

\[
t = \frac{\text{Coeficiente}}{\text{Error estándar}}
\]

Usando errores estándar usuales:
\begin{itemize}
       \item \(cumgpa\): \[
    t = \frac{0.193}{0.064} = 3.02 \quad (\text{Mayor que } 1.96 \rightarrow \text{Significativo})
    \]
    \item \(tothrs\): \[
    t = \frac{0.0014}{0.0012} = 1.17 \quad (\text{Menor que } 1.96 \rightarrow \text{No significativo})
    \]
\end{itemize}

 Usando errores estándar robustos:
\begin{itemize}
       \item \(cumgpa\): \[
    t = \frac{0.193}{0.074} = 2.61 \quad (\text{Mayor que } 1.96 \rightarrow \text{Significativo})
    \]
    \item \(tothrs\): \[
    t = \frac{0.0014}{0.0012} = 1.17 \quad (\text{Menor que } 1.96 \rightarrow \text{No significativo})
    \]
\end{itemize}

 Al nivel de significancia del 5\%,  \(cumgpa\) es estadísticamente significativo, mientras que \(tothrs\) no lo es. Esta conclusión es consistente tanto con los errores estándar usuales como con los robustos.
 El uso de errores estándar robustos no afecta las conclusiones en este caso, pero son útiles para garantizar resultados confiables ante heterocedasticidad.

}
% \item ¿Por qué es razonable la hipótesis $H_0:\beta_{crsgpa}=1$. Pruebe esta hipótesis contra la alternativa
% de dos colas al nivel de 5\%, usando los dos errores estándar. Describa sus conclusiones.


% {\color{blue}


% La hipótesis nula y alternativa son:
% \[
% H_0: \beta_{crsgpa} = 1 \]

% \[H_1: \beta_{crsgpa} \neq 1
% \]


% El estadístico \(t\) se calcula como:
% \[
% t = \frac{\widehat{\beta}_{crsgpa} - 1}{\text{Error estándar}}
% \]


% La estimación de \(\beta_{crsgpa}\) es \(0.9\). Los errores estándar son:
% \begin{itemize}
%     \item Usual: \(0.175\)
%     \item Robusto: \(0.166\)
% \end{itemize}

% \noindent El cálculo del estadístico \(t\) es el siguiente:

% \begin{itemize}
%     \item \textbf{Con errores estándar usuales:}
%     \[
%     t = \frac{0.9 - 1}{0.175} = \frac{-0.1}{0.175} = -0.57
%     \]
%     \item \textbf{Con errores estándar robustos:}
%     \[
%     t = \frac{0.9 - 1}{0.166} = \frac{-0.1}{0.166} = -0.60
%     \]
% \end{itemize}


% El valor absoluto del estadístico \(t\) es \(0.57\) (usual) y \(0.60\) (robusto), ambos menores que el valor crítico de \(1.96\). Por lo tanto, \textbf{no rechazamos} la hipótesis nula \(H_0: \beta_{crsgpa} = 1\) al nivel del 5\%.


%  Es razonable suponer que el coeficiente de \(crsgpa\) es igual a 1, ya que no hay evidencia suficiente para rechazar \(H_0\).
% Los resultados son consistentes independientemente del tipo de error estándar utilizado, lo que sugiere robustez en la conclusión frente a la heterocedasticidad.
% }

\item  Pruebe si el hecho de que el deporte del estudiante se practique en otoño tiene efecto sobre
el GPA del semestre, usando ambos errores estándar. ¿El nivel de significancia al que se
puede rechazar la prueba depende de cuál error estándar se emplee?

{\color{blue}
\[
H_0: \beta_{\text{season}} = 0 \quad \text{vs.} \quad H_1: \beta_{\text{season}} \neq 0
\]


El estadístico \(t\) se calcula como:
\[
t = \frac{\widehat{\beta}_{\text{season}}}{\text{Error estándar}}
\]


La estimación de \(\hat{\beta}_{\text{season}}\) es \(-0.157\). Los errores estándar son:
\begin{itemize}
    \item Usual: \(0.098\)
    \item Robusto: \(0.07\)
\end{itemize}

\noindent El cálculo del estadístico \(t\) es el siguiente:

\begin{itemize}
    \item \textbf{Con errores estándar usuales:}
    \[
    t = \frac{-0.157}{0.098} = -1.60
    \]
    \item \textbf{Con errores estándar robustos:}
    \[
    t = \frac{-0.157}{0.07} = -2.24
    \]
\end{itemize}


El valor crítico \(t_c\) para una prueba de dos colas al nivel del 5\% es \(1.96\):
\begin{itemize}
    \item Con errores estándar usuales, \(|t| = 1.60 < 1.96\). Por lo tanto, no rechazamos \(H_0\).
    \item Con errores estándar robustos, \(|t|=2.24 >1.96\). Se rechaza  \(H_0\).

    \item Usando errores estándar usuales, no hay suficiente evidencia para afirmar que practicar deportes en otoño afecta significativamente el \(trmgpa\).
    \item Usando errores estándar robustos, el efecto es  significativo al nivel del 5\%.
    \item La elección del tipo de error estándar afecta la interpretación. Si se sospecha de heterocedasticidad, se recomienda confiar en los errores robustos.
    \item El nivel de significancia al que se puede rechazar
la prueba depende de cu\'al error est\'andar se emplee.
\end{itemize}
}

\end{enumerate}
\end{document}